\documentclass{article}
\usepackage{amsmath}
\begin{document}
\section*{A typed bridge from broker evidence to old-duty consumption}
This development adapter separates three facts: an attempted send can
spend bytes; a named application callback can make a packet available;
an authoritative station may consume that packet to retry the same old
operation. None entails the next without a distinct premise. The
example uses the pinned broker fixture with a paid headroom read and
its previously checked civil warehouse station handoff.

\paragraph{Input and interpretation.}
An order-preserving local trace tap observes the native attempt
function's pre- and post-window rows, its returned admission value,
and the selected application callback. An attempt stores sender,
recipient, timestamp, offered bytes, the actual TX charge, and
radio admission. A callback is linked to its exact attempt occurrence
and checked for the authorized endpoint, size and sender. These
occurrence identifiers are local to a taped run, not authenticating
network identities. Within this synchronous fixture a radio-admitted
attempt must be followed by exactly one corresponding application
callback at the same time. That equality is an adapter assumption
checked here, not a property of an asynchronous deployment; in a
different runtime radio admission and later callback must be separate
events with an independent occurrence link.

The extractor replays just the relevant TX and RX windows by identity,
checks each transition against the raw tap, and preserves an attempt
that charges TX even when no callback arrives. An old row exactly one
second old is still live. No new RX debit is permitted on a rejected
attempt; late purging of old RX rows can depend on the position of
the rejection. Checks do not recover ambient RF truth from ledger
rows, infer that a tape is complete, or authenticate a fabricated
yet internally consistent tape. Common trust in this extracted
source path is a premise.

\paragraph{Application.}
At the station's cut read the local authority row, named permit and
snapshot tag in one read-only SQLite transaction. A consumer demands
the old keyed digest delivered by the authorized route and no later
than the chosen deadline, together with the matching fenced snapshot.
Only then may it install the receipt and retry. The native retry
supplies the cached result and the native old authority must still
be fenced. One radio result is backed by an equally sequenced earlier
copy taken before the fence; it gets the same packet at the same
time, but cannot accept it because its local tag and authority are
wrong. This negative uses different native station state, not an
invented warning bit or a renamed outcome. An earlier mission
deadline can reject the delivered packet without implying that
it was lost.

The callback-only projection also fails before the outcome is known.
The measured early callback lists are identical for shared and
independent report routing, including time, payload and logical
endpoint. In the former the report has already spent one byte in
the old receive ledger, so the old primary fails. In the latter
it spent that byte in the verifier ledger and the primary succeeds.
The attempt origin and actual resource identity supply a necessary
distinction for a feasibility statement at the cut. This is a
pair of source-generated cases, not a claim that every joined
deployment must need this tap.

Discarding the attempt debit is unsound for a later scheduling
decision. Treating callbacks alone as old authorization is also
unsound. Keeping full simulation state is not required by this
scoped consume function, but the local history, episode label
structure, receive route, named fence and timing cannot be silently
interchanged. This scoped tape demonstrates the evidence interface
that separate native source facts feed; it is not a generalized
runtime monitor, blind selection, an independently evaluated mission,
an exhaustive inventory disposition, or an advancement claim.
\end{document}
